\chapter{TMD 与全球拟合：与唯象学的对接}

格点纲领与物理世界其余部分之间有两个接口。纵向地，横向动量依赖分布
扩展了共线图像并带来自身的因子化物理——诞生于下文这篇论文。横向地，
全球拟合把成千上万的数据点变成 PDF 集合：格点计算既以之为基准，也
日益被其吸收为输入。本章收录 TMD 的奠基文献，加上两条拟合学派的
各一座地标。

% ----------------------------------------------------------------------
\section{背靠背喷注：TMD 物理的诞生}
\papercard{J.\ C.\ Collins, D.\ E.\ Soper}
{QCD 中的背靠背喷注}
{Nucl.\ Phys.\ B \textbf{193}, 381 (1981)}
{DOI: 10.1016/0550-3213(81)90506-7 | 姊妹篇：Collins--Soper--Sterman,
Nucl.\ Phys.\ B \textbf{250}, 199 (1985)}
\whyselected{横动量依赖因子化的定义性论文（显式约 $\sim$4/99，在库内
TMD 路线内普遍）：每个准 TMD 文章都建立在其概念之上。}
\physics{对横动量很小的色单态产生过程，共线因子化失效；观测量依赖一
类新的矩阵元——同时具有光锥与横向分离。Collins 与 Soper 引入快度切割
$\zeta$ 定义这些 TMD，证明其从硬、喷注与软扇区因子化，并推导快度中的
演化方程——Collins--Soper 核——辅以 $k_\perp$ 空间的 DGLAP：
\begin{equation}
  \frac{\partial \ln F(x,\mathbf{k}_\perp;\zeta,\mu)}
       {\partial \ln\sqrt{\zeta}} = K(b_\perp;\mu),
  \qquad b_\perp \leftrightarrow \mathbf{k}_\perp .
\end{equation}
姊妹篇 CSS 论文完成了 Drell--Yan 过程的证明。核 $K$ 在大 $b_\perp$ 处
内禀非微扰——这正是数十年后格点登场的位置。}
\impact{库内的准 TMD 波函数、软函数提取与非微扰 Collins--Soper 核确定
都以这个对象为目标；它们在格点上遭遇的快度发散重整化难题，正是此文
所发现结构的直接后代。}
\cardend

% ----------------------------------------------------------------------
\section{CT18：Hessian 学派的标准}
\papercard{T.-J.\ Hou et al.\ (CTEQ-TEA Collaboration)}
{利用 LHC 高精度数据的新 CTEQ 量子色动力学全局分析}
{Phys.\ Rev.\ D \textbf{103}, 014013 (2021)}
{arXiv:1912.10053 | DOI: 10.1103/PhysRevD.103.014013}
\whyselected{本库原文之一，也是其被引最多的唯象基准（拟合语料内约
$\sim$6/99）：语料中几乎每项格点 PDF 提取的默认对照集。}
\physics{CT18 把 CTEQ 方法论推进到完整的 LHC 时代：在扩展基上的
Hessian 误差分析、含重夸克方案变体的 NNLO 理论、以及对数据集张力
（尤其是 ATLAS 7 TeV W/Z 测量，由 CT18A/CT18Z 变体容纳）的审慎处理。
它交付标准烛光预言（Higgs 道、W 质量敏感性），供对撞机分析与第一性
原理输入共同校验。其前身 CT14（Dulat 等，PRD \textbf{93}, 033006
(2016)，亦为本库原文）确立参数化框架；Hessian 方法本身可追溯到
Pumplin 等，JHEP \textbf{07}, 012 (2002)。}
\impact{在本库内，CT18 以参考曲线身份出现于胶子 PDF、准 PDF 与分布
振幅的全部对比图中；"格点 PDF 进入全球拟合"的计划正是通过追问当格点
矩阵元加入数据池后 CT 类不确定度收缩多少来校准的。}
\cardend

% ----------------------------------------------------------------------
\section{NNPDF3.1：蒙特卡罗学派}
\papercard{R.\ D.\ Ball et al.\ (NNPDF Collaboration)}
{来自高精度对撞机数据的部分子分布}
{Eur.\ Phys.\ J.\ C \textbf{77}, 663 (2017)}
{arXiv:1706.00428 | DOI: 10.1140/epjc/s10052-017-5199-5}
\whyselected{CT18 的蒙特卡罗方法论对应版（约 $\sim$8/99），本库原文
之一；其后续 NNPDF4.0（《通往百分之一精度的质子结构》）同样在库内，
并以之为基线引用。}
\physics{NNPDF 用神经网络系综表示 PDF 并以蒙特卡罗拟合，用闭合检验
认证不确定度既不低估也不夸大真值。3.1 版新增独立粲 PDF（理论驱动但由
数据拟合）、top 微分分布与高精度 LHC 电弱玻色子数据，进入理论误差
（NNLO 劈裂函数、电弱修正）主导的精度区。其方法论链条——闭合检验、
副本压缩、正性与求和规则约束——成为领域的质量标准，并通过 NNPDF4.0
的超参数优化而收敛于机器学习工程本身。}
\impact{两个拟合学派共同框定了本语料的每项格点结果：库内各图同时叠加
CT18 与 NNPDF 频带。更深的汇流——全球拟合即机器学习、格点采样即生成
模型——正是下一章的主题。}
\cardend
